% ---------------------------------------------------

\section{Sailboat Physics}
\label{app:sailboat_physics}

\textbf{Lift Forces:} A sailboat is driven by two opposing lift forces. The sail generates lift from the wind, directed roughly perpendicular to the sail surface. The keel generates an opposing lateral lift force from the water. The net effect of these two forces propels the boat forward.

\textbf{Point of Sail:} The point of sail is the boat's angle relative to the wind direction. The angle of attack is the sail's angle relative to the apparent wind. The optimal angle of attack depends on the point of sail. The no-sail zone encompasses approximately $\pm45^{\circ}$ directly into the wind ($90^{\circ}$ total arc), where the sail cannot generate useful lift.

\textbf{Rudder Control:} The rudder modifies the interaction between sail and keel forces to change heading. Turning the rudder toward the wind decreases the effective keel force, causing the boat to turn away from the wind (bear away). Turning the rudder away from the wind increases the effective keel force, causing the boat to turn into the wind (head up).

\textbf{Center of Effort (COE) and Center of Lateral Resistance (CLR):} The COE is the point where the resultant sail lift force acts. The CLR is the point where the resultant keel/hull lateral resistance acts. When the COE is positioned slightly aft of the CLR, the boat exhibits weather helm---a natural tendency to turn into the wind. Weather helm is desirable because it provides inherent stability: if the helm is released, the boat rounds up into the wind and the sails luff, bringing the boat to a safe stop rather than accelerating uncontrolled.

% ---------------------------------------------------
